\begin{table}[tp]
\centering
\caption{Stochastic two-sided PL minimax optimization: \href{https://github.com/OptaResearch/auto-math/blob/main/LongTails/Setups/StochasticTwoSidedPL.lean}{Stochastic\allowbreak{}Two\allowbreak{}Sided\allowbreak{}PL}.}\label{tab:StochasticTwoSidedPL}
\parbox{\linewidth}{\footnotesize Real inner-product spaces $E,F$, with $E$ finite-dimensional and $F$ complete and second-countable, carry their Borel structures. Partial gradients are $L$-Lipschitz in $\|x-x'\|+\|y-y'\|$; $f(\cdot,y)$ and $-f(x,\cdot)$ satisfy PL inequalities against every competitor, with $0<\mu_x,\mu_y\le L$. Each $\max_y f(x,y)$ is attained. Jointly measurable, integrable, unbiased partial-gradient oracles have integrable squared errors and variance at most $\sigma^2$, with independent samples from a probability law. Known $\Delta\ge0$ satisfies $2\Phi(x_0)-\Phi^*-f(x_0,y_0)+\|\nabla_x f(x_0,y_0)\|^2/L\le\Delta$. Target: $\mathbb{E}[2\Phi(\widehat x)-\Phi^*-f(\widehat x,\widehat y)]\le\varepsilon$, $\varepsilon>0$.}\par\smallskip
\begin{threeparttable}
\footnotesize
\setlength{\tabcolsep}{4pt}\renewcommand{\arraystretch}{1.25}
\begin{tabularx}{\linewidth}{@{}>{\raggedright\arraybackslash}p{0.20\linewidth}>{\raggedright\arraybackslash}p{0.32\linewidth}>{\raggedright\arraybackslash}p{0.25\linewidth}>{\raggedright\arraybackslash}X@{}}
\toprule
\textbf{Result} & \textbf{Oracle complexity} & \textbf{Novelty} & \textbf{Limitations} \\
\midrule
\textbf{(1)} \href{https://github.com/OptaResearch/auto-math/blob/main/LongTails/Theorems/StochasticTwoSidedPLSmoothedAgda.lean}{Stochastic\allowbreak{}Two\allowbreak{}Sided\allowbreak{}PL\allowbreak{}Smoothed\allowbreak{}Agda} \textbf{(Ours)} \href{https://github.com/OptaResearch/auto-math/blob/main/scripts/reviews_after_review/StochasticTwoSidedPLSmoothedAgda.md}{(review)} & $\mathcal{O}\!\left(\left[\frac{L^2}{\mu_x\mu_y}+\frac{L^3\sigma^2}{\mu_x^2\mu_y^2\varepsilon}\right]\ell\right)$ & --- & --- \\
\midrule[1.2pt]
Stoc-AGDA \citep[Thm.~3.3]{stplYang2020} & $\mathcal{O}\!\left(1+\frac{L^3\Delta}{\mu_x\mu_y^2\varepsilon}+\frac{L^5\sigma^2}{\mu_x^2\mu_y^4\varepsilon}\right)$\newline Finite-dimensional $F$; saddle and slice minima attained. & \textbf{(1)}: Noise term smaller by $\kappa_y^2/\ell$ when $\kappa_y^2>\ell$; termwise comparison, not uniform dominance. & \textbf{(1)}: Random second-half iterate instead of last iterate; needs $\sigma,\Delta,\varepsilon$ additionally for tuning; $\Delta$ includes an initial gradient term. \\
\midrule
AGDA \citep[Thm.~3.2]{stplYang2020} & $\mathcal{O}\!\left(\frac{L^3}{\mu_x\mu_y^2}\ell\right)$\newline $\sigma=0$; finite-dimensional $F$; saddle and slice minima attained. & \textbf{(1)}: At $\sigma=0$, $\kappa_x\kappa_y$ instead of $\kappa_x\kappa_y^2$; strict condition-number gain as $\kappa_y$ grows. & \textbf{(1)}: Random second-half iterate instead of last iterate. \\
\midrule
Restarted Smoothed-AGDA; Derived bound \citep[Thm.~4.1]{stplYang2022} & $\mathcal{O}\!\left(\frac{L^2\ell}{\mu_x\mu_y}+\frac{L^3\sigma^2}{\mu_x^2\mu_y^2\varepsilon}\right)$\newline Finite-dimensional $F$; conditional on geometric restarts with deterministic bounds controlling expected stage initialization gaps. & i) \textbf{(1)}: Scope: Hilbert-space $F$.\newline ii) \textbf{(1)}: Single-run potential-PL analysis and formalization; matches deterministic rate, without restarts. & \textbf{(1)}: Extra $\ell$ on noise term if the expected-initialization restart extension holds; no established complexity gain over this conditional bound. \\
\midrule
FESS-GDA \citep[Thm.~3.5]{stplShen2024v2}; one client & Different target:\newline $\mathbb{E}\|(x_T,y_T)-(x_\star,y_\star)\|^2\le\delta^2$.\newline Finite-dimensional $F$; saddle and slice minima attained. Rates not comparable. & \textbf{(1)}: Scope: no saddle-existence assumption. & --- \\
\bottomrule
\end{tabularx}
\begin{tablenotes}[flushleft]\footnotesize
\item[] \textbf{Abbr:} $\Phi(x)=\max_y f(x,y)$; $\Phi^*=\inf_x\Phi(x)$; $\ell=1+\log(1+\Delta/\varepsilon)$; $\kappa_x=L/\mu_x$, $\kappa_y=L/\mu_y$; PL: Polyak--\L{}ojasiewicz. One call returns one stochastic partial gradient; costs include both branches and every restart stage. $\delta$ is distance accuracy; $(x_\star,y_\star)$ is a saddle point.
\end{tablenotes}
\end{threeparttable}
\end{table}
\clearpage
